
\begin{obeassessment}
\textbf{Kuis Bab 8}
1. Apa perbedaan utama antara struktur internal DES (Feistel) dan AES (Substitusi-Permutasi)?
2. Jelaskan peran \textit{S-box} dalam AES dan bagaimana ia dibentuk secara matematis.
3. Mengapa operasi dalam AES dilakukan dalam $\text{GF}(2^8)$ dan bukan aritmetika bilangan bulat biasa?
4. Tuliskan urutan transformasi pada satu putaran dekripsi AES dan jelaskan mengapa
   urutan itu tidak dapat diperoleh hanya dengan membalik urutan transformasi pada
   enkripsi.
5. Sebuah implementasi mandiri menghasilkan cipherteks yang berbeda dari
   implementasi lain untuk plainteks dan kunci yang sama. Sebutkan dua penyebab
   yang paling mungkin dan cara memeriksanya.
6. Diketahui state pada akhir putaran ke-9 sama dengan
   \[
   \code{EB40F21E592E38848BA113E71BC342D2},
   \]
   dan kunci putaran terakhir $K_{10}$ sama dengan
   \[
   \code{A150BE1DEFC16E55C4C018A35610A1DD}.
   \]
   Hitunglah cipherteksnya, dengan memperhatikan bahwa putaran terakhir tidak
   memuat MixColumns.
7. Bangkitkan $w[8]$ beserta $K_1$ untuk kunci AES-256 yang tersusun atas delapan
   word berikut:
   \[
   \code{603DEB10}\ \code{15CA71BE}\ \code{2B73AEF0}\ \code{857D7781}
   \]
   \[
   \code{1F352C07}\ \code{3B6108D7}\ \code{2D9810A3}\ \code{0914DFF4}.
   \]
   Jelaskan aturan mana yang Anda pakai, dan tunjukkan pada word keberapakah
   aturan khusus AES-256 mulai berlaku.
\end{obeassessment}
\begin{competencychecklist}
\begin{tabularx}{\linewidth}{|X|c|}
\hline
Indikator Kompetensi & Tercapai \\
\hline
Saya dapat menjelaskan alasan penggantian DES menjadi AES & [ ] \\
\hline
Saya dapat membedakan varian AES-128, 192, dan 256 & [ ] \\
\hline
Saya dapat menguraikan alur kerja \textit{SubBytes, ShiftRows, MixColumns} & [ ] \\
\hline
Saya dapat menjelaskan konsep matriks \textit{state} pada AES & [ ] \\
\hline
Saya dapat memahami dasar operasi aritmetika pada $\text{GF}(2^8)$ & [ ] \\
\hline
Saya dapat menjelaskan urutan transformasi pada dekripsi AES & [ ] \\
\hline
Saya dapat membedakan \textit{straightforward} dan \textit{equivalent inverse cipher} & [ ] \\
\hline
Saya dapat membangkitkan S-box AES dari konstruksi $\text{GF}(2^8)$ & [ ] \\
\hline
Saya dapat menghitung kunci putaran AES-128 dan AES-256 dengan tangan & [ ] \\
\hline
Saya dapat menunjukkan bahwa implementasi AES saya cocok dengan vektor uji FIPS 197 & [ ] \\
\hline
\end{tabularx}
\end{competencychecklist}
